\chapter{Governance aspects}
\label{chapter:governance}

An AI system that audits AI systems must itself be governed. We analyse the auditor along the course's governance
dimensions; Table~\ref{tab:agency} summarises its agency.

\section{Ethics}
\label{section:ethics}

The main ethical risk is \emph{false assurance}: a report that says ``compliant'' where the law is not met can harm the
people affected by the audited system (rejected candidates, declined borrowers) and shift blame to a tool. We therefore
(i) treat a false compliance as the worst error in both the metrics and the design (no positive verdict without a
verified citation; disagreement $\Rightarrow$ human); (ii) state in every report that it is decision support for an
internal audit~\cite{raji2020closing}, not a conformity assessment or legal advice; (iii) never let generated artefacts
invent facts (explicit \texttt{[TODO]} placeholders). \emph{Privacy:} the auditor reads source repositories that may
contain secrets or personal data in sample logs; only bounded excerpts are sent to the LLM, through the user's own Claude
account, with no tool access and no session persistence, and the local cache can be disabled or wiped. A deployment on
confidential code should add secret/PII redaction before prompting, which is listed as future work. \emph{Dual use:} the
same tool could be used to produce a ``clean'' report by tuning rules to pass; the full rule trace and the report history
make such manipulation visible.

\section{Fairness}
\label{section:fairness}

Fairness concerns the auditor in two ways. First, \emph{fairness of the audited systems}: the requirements include data
governance and bias examination (AI Act Art.~10), and the benchmark contains a system that uses age, gender and nationality
as features; the auditor can check that bias evaluation is \emph{documented and performed}, but it does not measure bias
itself (it has no access to data or predictions), and it should not be read as a fairness certification~\cite{mehrabi2021survey}.
Second, \emph{fairness of the auditor across the audited}: the same requirement should receive the same verdict for
equivalent evidence regardless of programming language, documentation language or naming style. Rules are biased toward
the vocabulary they were written for (e.g.\ English identifiers, \texttt{logger}); the LLM is more robust to naming but
may favour verbose, well-documented repositories. We report per-category and per-application results to expose such
disparities; extending the benchmark to non-English documentation and more stacks is future work.

\section{Agency: what the system may decide and how far it may act}
\label{section:agency}

Following the levels-of-automation view of Parasuraman et al.~\cite{parasuraman2000levels}, the auditor automates
information acquisition and analysis fully, \emph{recommends} decisions, and \emph{executes} nothing.

\begin{table}[htbp]
\caption{Agency of the auditor: decisions it may take alone, with a human, or never.}
\label{tab:agency}
\centering\small
\begin{tabularx}{\linewidth}{@{}p{3.3cm}XX@{}}
\toprule
\textbf{Level} & \textbf{Decisions} & \textbf{Condition} \\
\midrule
Autonomous & collect evidence, classify risk tier, mark \stat{na}, emit \stat{c}/\stat{p}/\stat{nc} verdicts,
draft remediation and artefacts, fail a CI build on a regression & verdict verified, confidence $\ge \beta$, no rule/LLM
conflict; CI failure is reversible and threshold set by the team \\
Human-in-the-loop & any verdict with confidence $< \beta$, any rule vs.\ LLM conflict, any \stat{c} claim without valid
citation, contested verdicts & routed to the review queue as \stat{nr}; a named reviewer decides \\
Never & modify the audited repository, commit or merge, file anything with a regulator, declare legal conformity,
decide applicability by LLM alone & out of scope by design (no write tools; judge cannot output \stat{na}) \\
\bottomrule
\end{tabularx}
\end{table}

The system may therefore act \emph{only on its own outputs} (a report, a SARIF file, a CI exit code). The furthest it acts
before a person is involved is blocking a merge in CI, which is reversible, configurable (fail-on severity) and based on a
regression relative to a baseline that a human accepted.

\section{Architecture: components required by that permission}
\label{section:architecture-gov}

Each permission in Table~\ref{tab:agency} is backed by a concrete component:
\begin{itemize}[leftmargin=*]
  \item \textbf{Retrieval}: the evidence index and the BM25 EvidenceScout bound what the judge can see; every claim must be
        traceable to an indexed chunk with path and line span.
  \item \textbf{Tools}: none for the LLM. The Claude Code CLI runs with tools, MCP servers and settings disabled in an empty
        temporary directory; only the deterministic code reads files. This removes the risk of an agent modifying the repo.
  \item \textbf{Memory}: no conversational memory between calls (no session persistence); the only memory is the
        content-addressed response cache (reproducibility) and the stored reports with their review history
        (institutional memory, input to CI regression checks).
  \item \textbf{Orchestration}: a fixed, ordinary-code pipeline (Algorithm~\ref{alg:hybrid}) with explicit routing
        thresholds; agents cannot spawn agents or change the plan.
  \item \textbf{Checks}: JSON-schema validation with bounded retries, the Verifier (Algorithm~\ref{alg:verifier}), the
        conflict rule and the confidence gate, and the CI regression comparator.
\end{itemize}

\section{Accountability: who answers, what the record shows, how to contest}
\label{section:accountability}

Accountability for AI requires an identifiable answerable agent, a record and a forum~\cite{novelli2023accountability}.
\begin{itemize}[leftmargin=*]
  \item \textbf{Who answers.} The organisation deploying the audited system (the AI Act provider or deployer) remains
        responsible for compliance; within it, the named reviewer who resolves an item answers for that decision, and the
        maintainers of the requirements YAML answer for the rules. The tool's authors answer for the correctness of the
        pipeline and its documented limits.
  \item \textbf{What the record shows.} Each report stores the policy ids, the full requirements used, the mode,
        model and backend, the risk tier with triggers, every verdict with method, confidence, verified flag, rationale,
        cited evidence (path, lines, snippet), gaps and remediation, plus timings, LLM calls, tokens and cost. Human
        overrides append ``\emph{Human review by X: old $\to$ new. Justification}'' and are listed in the report statistics (\texttt{human\_reviews});
        the cache makes LLM outputs replayable.
  \item \textbf{How to contest.} Any verdict can be contested in the dashboard review queue or via
        \texttt{ccauditor review} / \texttt{POST /reports/review}; the contester sees the exact evidence and both
        rationales, and the decision is recorded with name and justification, in line with contestability by
        design~\cite{alfrink2022contestable}. Rule-level errors are contested by changing the versioned requirement YAML,
        and the policy-diff tool shows which reports are affected by such a change.
\end{itemize}

\paragraph{Transparency.} The rules, prompts, thresholds and benchmark are open in the repository; the report exposes the
method that produced each verdict, so a reader can distinguish rule results, LLM judgements and human decisions.
